% =============================================================================
% 11-trees-bst-avl-notes.tex — Lecture Notes for Week 11:
% Keeping Order (Trees, Binary Search Trees & AVL)
%
% DERIVED ARTIFACT. Source of truth:
% Slides/CS301/11-trees-bst-avl.tex (60 pp incl. title/transitions).
% Generated by /lecture-notes at the "textbook-candidate prose" Detail Bar
% (every trace step written out, every example fully solved, Exercises with a
% separate Solutions block, reading guidance per section); do not hand-edit
% content independently of the Beamer source without re-running /qa-notes
% afterward (single-source-of-truth.md).
%
% Page-grounding (inherited 1:1 from the deck header; no upgrades):
%   Wirth2004 Sec.4.4 Tree Structures p.126, Sec.4.4.2 Basic Operations
%   p.132, Sec.4.4.3 Tree Search and Insertion p.135, Sec.4.4.4 Tree Deletion
%   p.139, Sec.4.5 Balanced Trees p.142;
%   Sedgewick2011 Sec.3.2 Binary Search Trees p.396, Sec.3.3 Balanced Search
%   Trees p.424 (motivation only -- its balanced tree is red-black, not AVL);
%   Karumanchi2017 Ch.6 Trees, pp.224-368 PDF.
%   Horowitz & Sahni Ch.5 is cited at chapter level only -- no page numbers
%   are asserted.
%   All TikZ diagrams are reused verbatim from the deck.
%
% Compile with (macOS/Linux, ':'-joined TEXINPUTS/BIBINPUTS): cd Notes/CS301 &&
%   TEXINPUTS=../../Preambles:$TEXINPUTS xelatex -interaction=nonstopmode 11-trees-bst-avl-notes.tex
%   BIBINPUTS=../..:$BIBINPUTS bibtex 11-trees-bst-avl-notes
%   TEXINPUTS=../../Preambles:$TEXINPUTS xelatex -interaction=nonstopmode 11-trees-bst-avl-notes.tex  (x2)
% =============================================================================

\documentclass[11pt]{article}
\usepackage[margin=1in]{geometry}
% =============================================================================
% Preambles/header.tex — shared LaTeX/Beamer preamble for this project
%
% Use from a Beamer lecture by prepending:
%     \input{header}
% and compiling with TEXINPUTS=../Preambles:$TEXINPUTS (the /compile-latex
% skill does this automatically).
%
% PALETTE CONTRACT. Color names defined here MUST match the names in
% Quarto/theme-template.scss so Beamer and Quarto renderings use the same
% palette. The scripts/check-palette-sync.sh script enforces this.
%
% To customize: edit the HEX values below AND the matching $variable in
% Quarto/theme-template.scss, then run scripts/check-palette-sync.sh.
% =============================================================================

% -----------------------------------------------------------------------------
% Engine detection + encoding
%
% Our /compile-latex skill uses XeLaTeX, where inputenc is unnecessary and
% can produce encoding warnings. Only load inputenc under pdfTeX.
% -----------------------------------------------------------------------------
\usepackage{iftex}
\ifPDFTeX
  \usepackage[utf8]{inputenc}
\fi

\usepackage{amsmath, amssymb, amsthm}
\usepackage{booktabs}
\usepackage{graphicx}

% -----------------------------------------------------------------------------
% PALETTE — mirrors Quarto/theme-template.scss variable names.
% Load xcolor and define colors BEFORE anything that references them
% (notably hyperref, hypersetup, and the Beamer color assignments).
% -----------------------------------------------------------------------------
\usepackage{xcolor}

\definecolor{primary-blue}{HTML}{012169}      % headings, accents
\definecolor{primary-gold}{HTML}{B9975B}      % emphasis, borders
\definecolor{highlight-yellow}{HTML}{F2A900}  % markers, alerts
\definecolor{light-bg}{HTML}{E8EDF5}          % title / section backgrounds
\definecolor{jet}{HTML}{1A1A1A}               % body text

% Semantic colors (used in diagrams and annotations)
\definecolor{positive}{HTML}{15803D}          % good / observed / identified
\definecolor{negative}{HTML}{B91C1C}          % bad / problematic / bias
\definecolor{neutral}{HTML}{525252}           % reference / context

% Highlight accents (match .hi-* classes in the SCSS)
\definecolor{hi-slate}{HTML}{314F4F}
\definecolor{hi-green}{HTML}{2E8B57}
\definecolor{hi-red}{HTML}{C41E3A}

% Semantic aliases so skill templates and snippets can write
% \textcolor{accent}{...} without hard-coding "primary-blue".
\colorlet{accent}{primary-blue}
\colorlet{accent2}{primary-gold}

% -----------------------------------------------------------------------------
% Hyperref — loaded AFTER xcolor + palette so link colors resolve.
% -----------------------------------------------------------------------------
\usepackage{hyperref}
\hypersetup{colorlinks=true, linkcolor=primary-blue, urlcolor=primary-blue,
            citecolor=primary-blue, pdfborder={0 0 0}}

% -----------------------------------------------------------------------------
% Course code — set once per deck (after \input{header}, before \begin{document})
% via \coursecode{CS401}. Shown in the footer on every slide so a deck's course
% is identifiable without opening the title page. Empty by default.
% -----------------------------------------------------------------------------
\makeatletter
\newcommand{\coursecode}[1]{\gdef\@coursecodetext{#1}}
\gdef\@coursecodetext{}
\makeatother

% -----------------------------------------------------------------------------
% Beamer theme assignments (only apply under Beamer).
%
% We detect Beamer via \@ifclassloaded{beamer}, which is the documented,
% non-internal way. This must be wrapped in \makeatletter/\makeatother
% because \@ifclassloaded contains an @-catcode character.
% -----------------------------------------------------------------------------
\makeatletter
\@ifclassloaded{beamer}{%
  \usefonttheme{professionalfonts}
  \setbeamercolor{structure}{fg=primary-blue}
  \setbeamercolor{frametitle}{fg=primary-blue}
  \setbeamercolor{title}{fg=primary-blue}
  \setbeamercolor{subtitle}{fg=primary-gold}
  \setbeamercolor{author}{fg=primary-blue}
  \setbeamercolor{institute}{fg=neutral}
  \setbeamercolor{date}{fg=primary-gold}
  \setbeamercolor{itemize item}{fg=primary-blue}
  \setbeamercolor{itemize subitem}{fg=primary-gold}
  \setbeamercolor{alerted text}{fg=primary-gold}
  \setbeamercolor{block title}{fg=white, bg=primary-blue}
  \setbeamercolor{block body}{bg=light-bg}
  % Semantic block variants: block=definition/concept (blue), exampleblock=
  % worked example (green/positive), alertblock=key takeaway or caution (gold).
  % Keep to <=2 colored boxes per slide across ALL three variants (INV-7).
  \setbeamercolor{block title example}{fg=white, bg=positive}
  \setbeamercolor{block body example}{bg=light-bg}
  \setbeamercolor{block title alerted}{fg=white, bg=primary-gold}
  \setbeamercolor{block body alerted}{bg=light-bg}

  % Minimal footer: course code (left) + page X / N (right), no navigation clutter
  \setbeamertemplate{navigation symbols}{}
  \setbeamertemplate{footline}{%
    \color{neutral}\footnotesize\hspace{0.7em}\@coursecodetext\hfill
    \insertframenumber/\inserttotalframenumber\hspace{0.7em}\vspace{0.3em}}
}{}
\makeatother

% -----------------------------------------------------------------------------
% TikZ setup — libraries the snippet gallery uses
% -----------------------------------------------------------------------------
\usepackage{tikz}
\usetikzlibrary{arrows.meta, positioning, calc, decorations.pathreplacing,
                fit, shapes.geometric, backgrounds}

% Shared TikZ styles — snippets and hand-written diagrams can pick these up
% via `\begin{tikzpicture}[every node/.style={...}]` or by naming specific
% styles (e.g., `\node[dag-node] (X) at (0,0) {X};`).
\tikzset{
  % Node styles for causal DAGs and similar
  dag-node/.style = {
    draw=primary-blue, fill=light-bg, rounded corners=2pt,
    minimum width=1.2cm, minimum height=0.9cm, align=center, font=\small
  },
  decision-node/.style = {
    draw=primary-gold, fill=white, diamond, aspect=1.8,
    minimum width=1.8cm, minimum height=1.2cm, align=center, font=\small,
    inner sep=1pt
  },
  % Edge styles — observed vs counterfactual
  observed-edge/.style = {-{Latex[length=2.5mm]}, thick, draw=jet},
  counterfactual-edge/.style = {-{Latex[length=2.5mm]}, thick, draw=neutral, dashed},
  confound-edge/.style = {-{Latex[length=2.5mm]}, thick, draw=negative, dashed},
  % Data-point styles for plots
  observed-dot/.style = {fill=primary-blue, draw=primary-blue, circle, inner sep=1.2pt},
  counterfactual-dot/.style = {fill=white, draw=neutral, circle, inner sep=1.2pt}
}

% -----------------------------------------------------------------------------
% Convenience macros
% -----------------------------------------------------------------------------
\newcommand{\muted}[1]{\textcolor{neutral}{#1}}
\newcommand{\key}[1]{\textcolor{primary-gold}{\textbf{#1}}}
\newcommand{\good}[1]{\textcolor{positive}{#1}}
\newcommand{\bad}[1]{\textcolor{negative}{#1}}

% Standout frame macro: full-bleed dark background for section transitions.
% Beamer-only (uses \begin{frame}/\setbeamercolor/\usebeamerfont) -- do not
% call from an article-class file (e.g. Notes/<CODE>/*-notes.tex). \muted,
% \key, \good, \bad, and \coursecode above are plain xcolor/gdef and are
% safe to use from either class; verified when Notes/article-class support
% was added.
% Expand: \transitionslide{Your transition title}
\newcommand{\transitionslide}[1]{%
  \begingroup
  \setbeamercolor{background canvas}{bg=primary-blue}
  \begin{frame}[plain]
    \centering
    \vfill
    {\usebeamerfont{title}\color{white}\Large #1\par}
    \vfill
  \end{frame}
  \endgroup
}

% Section divider frame (Beamer-only), an alternative to \transitionslide with a
% darker two-line style: near-black (jet) background, a small white label line
% above a larger gold title, and a thin gold rule along the bottom edge.
% Expand: \sectiondivider{Section 2}{Pointers}
\newcommand{\sectiondivider}[2]{%
  \begingroup
  \setbeamercolor{background canvas}{bg=jet}
  \begin{frame}[plain]
    \vspace*{3.0cm}
    {\color{white}\ttfamily\large #1\par}
    \vspace{0.5cm}
    {\color{primary-gold}\LARGE #2\par}
    \begin{tikzpicture}[remember picture, overlay]
      \draw[primary-gold, line width=2.5pt]
        ($(current page.south west)+(0,0.1)$) -- ($(current page.south east)+(0,0.1)$);
    \end{tikzpicture}
  \end{frame}
  \endgroup
}

% -----------------------------------------------------------------------------
% Channel watermark — "Concepts That Click" (YouTube).
%
% Article-class files (CompetitiveExam Q/A sets, Notes, Assignments) get a
% light diagonal draft watermark on every page plus a centered footer naming
% the channel. Beamer decks get a subtle TikZ background watermark instead
% (the footline already carries the course code). Centralized here so every
% MCQ PDF is branded without per-file edits.
% -----------------------------------------------------------------------------
\makeatletter
\@ifclassloaded{article}{%
  \usepackage{draftwatermark}
  \SetWatermarkText{Concepts That Click}
  \SetWatermarkScale{1}
  \SetWatermarkFontSize{2cm}
  \SetWatermarkLightness{0.86}
  \SetWatermarkAngle{45}
  \usepackage{fancyhdr}
  \pagestyle{fancy}
  \fancyhf{}
  \fancyfoot[C]{\footnotesize\textcolor{neutral}{Concepts That Click~|~YouTube: \href{https://www.youtube.com/@concepts.thatclick}{@concepts.thatclick}}}
  \renewcommand{\headrulewidth}{0pt}
  \renewcommand{\footrulewidth}{0pt}
  \fancypagestyle{plain}{%
    \fancyhf{}%
    \fancyfoot[C]{\footnotesize\textcolor{neutral}{Concepts That Click~|~YouTube: \href{https://www.youtube.com/@concepts.thatclick}{@concepts.thatclick}}}%
    \renewcommand{\headrulewidth}{0pt}%
    \renewcommand{\footrulewidth}{0pt}%
  }
}{}
\@ifclassloaded{beamer}{%
  \setbeamertemplate{background}{%
    \begin{tikzpicture}[remember picture, overlay]
      \node[rotate=45, opacity=0.055, align=center]
        at (current page.center)
        {\Huge\textcolor{neutral}{Concepts That Click}};
    \end{tikzpicture}%
  }
}{}
\makeatother 
\coursecode{CS301}

% Emergency stretch so long \texttt{} tokens in prose do not spill the margin.
\setlength{\emergencystretch}{2.5em}

% Chapter-style numbering: this is Week 11's chapter.
\renewcommand{\thesection}{11.\arabic{section}}
\renewcommand{\thefigure}{11.\arabic{figure}}
\renewcommand{\thetable}{11.\arabic{table}}
\theoremstyle{definition}
\newtheorem{example}{Example}[section]

\newenvironment{definitionbox}[1]{%
  \par\noindent\textbf{Definition (#1).}\ \itshape
}{\par\vspace{0.3em}}

% Same tree styles as the Beamer deck, so figures are verbatim ports.
\tikzset{
  tnode/.style   = {circle, draw=primary-blue, fill=light-bg,
                    minimum size=0.80cm, inner sep=0pt, font=\small},
  tnodeg/.style  = {circle, draw=primary-gold, line width=1.1pt,
                    fill=highlight-yellow!25, minimum size=0.80cm,
                    inner sep=0pt, font=\small},
  tnodesm/.style = {circle, draw=primary-blue, fill=light-bg,
                    minimum size=0.55cm, inner sep=0pt, font=\scriptsize},
  tnodesmg/.style= {circle, draw=primary-gold, line width=1pt,
                    fill=highlight-yellow!25, minimum size=0.55cm,
                    inner sep=0pt, font=\scriptsize},
  tedge/.style   = {draw=neutral, thick}
}

\title{Keeping Order\\\large Lecture Notes --- Week 11: Trees, Binary Search
  Trees \& AVL}
\author{Auqib Hamid Lone\\[0.3em]
  \small Department of Computer Science \& Engineering\\
  \small Government College of Engineering \& Technology (GCET), Ganderbal}
\date{\today}

\begin{document}
\maketitle

\section{Why This Week Exists: Order Is the Missing Operation}

\textbf{Read alongside:} Wirth, \textit{Algorithms and Data Structures},
Sec.~4.4, pp.126--149 (Oberon revision; cite the section number if you hold
the 1985 print edition); Sedgewick and Wayne, \textit{Algorithms}, Sec.~3.2,
p.396; Karumanchi, \textit{Data Structures and Algorithms Made Easy}, Ch.~6,
pp.224--368 (PDF). Horowitz and Sahni, Ch.~5, is the institute's prescribed
chapter, cited here at chapter level only
\cite[Ch.~5]{HorowitzSahni2008_fundamentals_data_structures}.

Week~10 ended with a fast but one-sided structure. A hash table finds a
record in about one probe, but it cannot answer ``give me the next key'' or
``give me every key between $a$ and $b$'', because it deliberately throws
order away. The closing sentence of Week~10 was a promise: \emph{Week 11
builds trees --- binary trees, BSTs and AVL rotations --- that keep order and
logarithmic search.} This chapter pays that promise off.

The deficiency is easiest to see as a table. For the calculator's symbol
table of $2{,}000$ identifiers, the three structures we already have each
give up something:

\begin{center}
{\footnotesize
\begin{tabular}{l c c c}
  \toprule
  \textbf{structure} & \textbf{search} & \textbf{insert} & \textbf{ordered?} \\
  \midrule
  sorted array (Weeks 8--9) & $O(\log n)$ & $O(n)$ & yes \\
  hash table (Week 10)      & $O(1)$ expected & $O(1)$ expected & no \\
  linked list (Week 6)      & $O(n)$ & $O(1)$ at head & no \\
  \midrule
  balanced BST (this week)  & $O(\log n)$ & $O(\log n)$ & yes \\
  \bottomrule
\end{tabular}}
\end{center}

We want the last row: fast search, fast insertion, and order. A tree is the
first structure in this course that can give all three at once. The price is
that we must control its \emph{shape}: the same keys can make a short tree or
a long chain, and the cost follows the shape.

The running application remains the calculator. Week~10 could find an
identifier quickly but could not print the identifiers in alphabetical order.
This week's tree can do both, and it can answer range questions such as
``which identifiers begin with \texttt{ab}?'' without scanning the table.

\section{Tree Vocabulary}

\subsection{Nodes, edges, and the words we use}

\begin{definitionbox}{tree}
A \key{tree} is a set of \key{nodes} joined by \key{edges}, with one
distinguished node called the \key{root}, such that there is exactly one path
from the root to every node and no cycles. A node directly below another is
its \key{child}; the node above is its \key{parent}. Nodes with the same
parent are \key{siblings}. A \key{leaf} is a node with no children.
\end{definitionbox}

Figure~\ref{fig:anatomy} shows the seven-node tree used throughout this
chapter. The root is 50. The leaves are 20, 40, 60 and 80. The parent of 20
is 30, the child of 30 is 20, and 20 and 40 are siblings. A \key{subtree} is
any node together with everything below it: 30, 20 and 40 form the left
subtree of 50.

\begin{figure}[h]
  \centering
  \begin{tikzpicture}[every node/.style={font=\small}]
    \node[tnode] (n50) at (0,3)   {50};
    \node[tnode] (n30) at (-1.7,1.8) {30};
    \node[tnode] (n70) at (1.7,1.8)  {70};
    \node[tnode] (n20) at (-2.5,0.6) {20};
    \node[tnode] (n40) at (-0.9,0.6) {40};
    \node[tnode] (n60) at (0.9,0.6)  {60};
    \node[tnode] (n80) at (2.5,0.6)  {80};
    \draw[tedge] (n50) -- (n30);
    \draw[tedge] (n50) -- (n70);
    \draw[tedge] (n30) -- (n20);
    \draw[tedge] (n30) -- (n40);
    \draw[tedge] (n70) -- (n60);
    \draw[tedge] (n70) -- (n80);
    \foreach \y/\t in {3/0, 1.8/1, 0.6/2}
      \node[text=neutral, anchor=east, font=\scriptsize] at (-3.7,\y)
        {depth \t};
  \end{tikzpicture}
  \caption{The running seven-node tree. Depth is counted from the root; the
    tree has height 2. Cf.\ Wirth, Sec.~4.4, pp.126--131
    \cite{Wirth2004_algorithms_data_structures}.}
  \label{fig:anatomy}
\end{figure}

\subsection{Depth and height}

\begin{definitionbox}{depth, height}
The \key{depth} of a node is the number of edges from the root down to that
node; the root has depth $0$. The \key{height} of a node is the number of
edges on the longest path from that node down to a leaf; a leaf has height
$0$. The \key{height of a tree} is the height of its root.
\end{definitionbox}

For Figure~\ref{fig:anatomy}, node 50 has depth $0$ and height $2$; nodes 30
and 70 have depth $1$ and height $1$; the four leaves have depth $2$ and
height $0$.

\textbf{A convention must be stated.} In this course a single node has height
$0$ and the empty tree has height $-1$. Some books count nodes instead of
edges, so the same tree has height $3$ under that convention. An exam answer
that says only ``height 3'' without saying which convention it uses is
ambiguous; write the convention down.

\subsection{Binary trees and three useful shapes}

\begin{definitionbox}{binary tree}
A \key{binary tree} is a tree in which every node has at most two children,
called its \key{left child} and its \key{right child}. The two sides are
ordered: swapping a left and a right child produces a different binary tree,
even when the drawing looks similar.
\end{definitionbox}

In C the node is the Week~6 linked-list node with two links instead of one:

\begin{verbatim}
struct node {
    int data;
    struct node *left;
    struct node *right;
};
\end{verbatim}

Three shapes are worth naming because textbooks and exam questions use them:

\begin{itemize}
  \item \key{Full}: every node has either $0$ or $2$ children. The seven-node
    tree of Figure~\ref{fig:anatomy} is full.
  \item \key{Complete}: every level is full except possibly the last, and the
    last level is filled from left to right.
  \item \key{Perfect}: the tree is full and all leaves are at the same depth.
    The seven-node tree is perfect.
\end{itemize}

\begin{figure}[h]
  \centering
  \begin{tikzpicture}[every node/.style={font=\scriptsize}]
    \foreach \x/\y/\v in {1/2/4, 0.4/1/2, 1.6/1/6,
                          0.1/0/1, 0.7/0/3, 1.3/0/5, 1.9/0/7}
      \node[tnodesm] at (\x,\y) {\v};
    \draw[tedge] (1,2) -- (0.4,1);
    \draw[tedge] (1,2) -- (1.6,1);
    \draw[tedge] (0.4,1) -- (0.1,0);
    \draw[tedge] (0.4,1) -- (0.7,0);
    \draw[tedge] (1.6,1) -- (1.3,0);
    \draw[tedge] (1.6,1) -- (1.9,0);
    \node[text=neutral, font=\tiny] at (1,-0.75)
      {perfect: 7 nodes, leaves at depth 2};

    \foreach \x/\y/\v in {5.5/2/4, 4.9/1/2, 6.1/1/6,
                          4.6/0/1, 5.2/0/3, 5.8/0/5}
      \node[tnodesm] at (\x,\y) {\v};
    \draw[tedge] (5.5,2) -- (4.9,1);
    \draw[tedge] (5.5,2) -- (6.1,1);
    \draw[tedge] (4.9,1) -- (4.6,0);
    \draw[tedge] (4.9,1) -- (5.2,0);
    \draw[tedge] (6.1,1) -- (5.8,0);
    \node[text=neutral, font=\tiny] at (5.5,-0.75)
      {complete: last level fills left to right};
  \end{tikzpicture}
  \caption{A perfect tree (left) and a complete tree that is not perfect
    (right). Both are full. Shape names matter for the height bounds that
    follow.}
  \label{fig:shapes}
\end{figure}

\subsection{Counting nodes, and why a logarithm appears}

Four facts follow directly from the definition of a binary tree.

\begin{enumerate}
  \item Level $d$ holds at most $2^d$ nodes, with the root at level $0$:
    $1, 2, 4, 8, \dots$
  \item A tree of height $h$ holds at most
    \[
      1 + 2 + 4 + \cdots + 2^h = 2^{h+1} - 1
    \]
    nodes.
  \item A tree with $n$ nodes has height at least $\lfloor \log_2 n\rfloor$.
    A perfect tree with $n$ nodes has height exactly $\log_2(n+1)-1$.
  \item In any binary tree with $n$ nodes there are exactly $n+1$ empty child
    slots (\texttt{NULL} pointers), and a full binary tree with $L$ leaves
    has $L-1$ internal nodes.
\end{enumerate}

The fourth fact is worth deriving once, because it is a common exam
question. A binary tree has $2n$ child slots in total (two per node). Exactly
$n-1$ of them are occupied, because every node except the root is the child
of exactly one edge. Therefore the number of empty slots is
$2n-(n-1)=n+1$.

The logarithm in the third fact is the reason trees matter. Each level can
double the number of nodes, so a tree with $1{,}023$ nodes can be only $9$
edges tall. A search that walks from the root to a leaf then costs about $9$
comparisons, not $1{,}023$. Everything in the rest of this chapter depends on
keeping that short shape.

\section{Four Ways to Visit Every Node}

A \key{traversal} visits every node exactly once and prints its key in some
order. The order is not arbitrary: it is chosen by \emph{when} we print a
node relative to its two subtrees. Figure~\ref{fig:walk} is the tree we will
walk; all four traversals cost $O(n)$ time because each node is handled once.

\begin{figure}[h]
  \centering
  \begin{tikzpicture}[every node/.style={font=\small}]
    \node[tnode] (n50) at (0,2.8)   {50};
    \node[tnode] (n30) at (-1.7,1.7) {30};
    \node[tnode] (n70) at (1.7,1.7)  {70};
    \node[tnode] (n20) at (-2.5,0.6) {20};
    \node[tnode] (n40) at (-0.9,0.6) {40};
    \node[tnode] (n60) at (0.9,0.6)  {60};
    \node[tnode] (n80) at (2.5,0.6)  {80};
    \draw[tedge] (n50) -- (n30);
    \draw[tedge] (n50) -- (n70);
    \draw[tedge] (n30) -- (n20);
    \draw[tedge] (n30) -- (n40);
    \draw[tedge] (n70) -- (n60);
    \draw[tedge] (n70) -- (n80);
  \end{tikzpicture}
  \caption{The tree used for all four traversal traces. Cf.\ Wirth,
    Sec.~4.4.2, pp.132--135 \cite{Wirth2004_algorithms_data_structures};
    Sedgewick, Sec.~3.2, pp.396--423 \cite{Sedgewick2011_algorithms}.}
  \label{fig:walk}
\end{figure}

\subsection{Inorder: left, node, right}

\begin{definitionbox}{inorder traversal}
Traverse the left subtree, print the node, then traverse the right subtree.
\end{definitionbox}

Start at 50 and go left to 30, then left to 20. Node 20 has no left child, so
print \textbf{20}; it has no right child, so return to 30. Print \textbf{30},
then go right to 40: print \textbf{40}, return to 50. Print \textbf{50}, then
repeat on the right subtree: print \textbf{60}, return to 70, print
\textbf{70}, then \textbf{80}. The output is
\[
  20,\ 30,\ 40,\ 50,\ 60,\ 70,\ 80.
\]
That is sorted. It is not a coincidence: in a binary search tree, inorder
visits keys from smallest to largest. On an arbitrary binary tree, inorder is
merely ``left, node, right''.

\subsection{Preorder: node, left, right}

\begin{definitionbox}{preorder traversal}
Print the node, then traverse the left subtree, then the right subtree.
\end{definitionbox}

Print \textbf{50}; go left and print \textbf{30}; print \textbf{20}; return
and print \textbf{40}. Return to 50, go right, and print \textbf{70},
\textbf{60}, \textbf{80}:
\[
  50,\ 30,\ 20,\ 40,\ 70,\ 60,\ 80.
\]
Preorder is useful for copying a tree or writing it to a file: a parent is
always printed before its children, so the same shape can be rebuilt. It is
also the prefix form of an expression, as Section~11.4 shows.

\subsection{Postorder: left, right, node}

\begin{definitionbox}{postorder traversal}
Traverse the left subtree, traverse the right subtree, then print the node.
\end{definitionbox}

The left subtree of 50 produces 20, 40, then \textbf{30}. The right subtree
produces 60, 80, then \textbf{70}. Finally print the root \textbf{50}:
\[
  20,\ 40,\ 30,\ 60,\ 80,\ 70,\ 50.
\]
Postorder is the traversal to use when freeing a tree: children are freed
before their parent, so no live pointer is left dangling. It is also how an
expression tree is evaluated.

\subsection{Level order: top to bottom, left to right}

\begin{definitionbox}{level-order traversal}
Visit the root, then all nodes at depth $1$, then all at depth $2$, and so
on; within a level, go from left to right.
\end{definitionbox}

For Figure~\ref{fig:walk}: depth $0$ gives \textbf{50}; depth $1$ gives
\textbf{30, 70}; depth $2$ gives \textbf{20, 40, 60, 80}. The output is
\[
  50,\ 30,\ 70,\ 20,\ 40,\ 60,\ 80.
\]
This is not recursion on the two children. Use the queue from Week~5:
enqueue the root; while the queue is non-empty, dequeue a node, print it, and
enqueue its left and right children. Level order is breadth-first search on a
tree, and the same idea returns for graphs in Week~12.

Table~\ref{tab:orders} collects the four outputs.

\begin{table}[h]
  \centering
  {\footnotesize
  \begin{tabular}{l l}
    \toprule
    \textbf{traversal} & \textbf{output for the seven-node tree} \\
    \midrule
    inorder   & 20, 30, 40, 50, 60, 70, 80 \\
    preorder  & 50, 30, 20, 40, 70, 60, 80 \\
    postorder & 20, 40, 30, 60, 80, 70, 50 \\
    level order & 50, 30, 70, 20, 40, 60, 80 \\
    \bottomrule
  \end{tabular}}
  \caption{The same tree, four orders. Only inorder is sorted here because
    this particular tree is a BST.}
  \label{tab:orders}
\end{table}

A memory hook: \key{in}order puts the node \emph{in} the middle, \key{pre}order
puts it before, and \key{post}order puts it after.

\subsection{Traversals in C}

The three depth-first traversals differ only in where the \texttt{printf}
sits:

\begin{verbatim}
void inorder(struct node *t) {
    if (t == NULL) return;
    inorder(t->left);
    printf("%d ", t->data);
    inorder(t->right);
}

void preorder(struct node *t) {
    if (t == NULL) return;
    printf("%d ", t->data);
    preorder(t->left);
    preorder(t->right);
}

void postorder(struct node *t) {
    if (t == NULL) return;
    postorder(t->left);
    postorder(t->right);
    printf("%d ", t->data);
}
\end{verbatim}

Each function visits every node once and does $O(1)$ work per node, so each
is $O(n)$. The recursion stack is $O(h)$ in the worst case, where $h$ is the
height of the tree.

\begin{example}
Walk the tree of Figure~\ref{fig:check} and write its inorder, preorder and
postorder outputs.

\begin{figure}[h]
  \centering
  \begin{tikzpicture}[every node/.style={font=\scriptsize}]
    \node[tnodesm] (a) at (0,2.6) {8};
    \node[tnodesm] (b) at (-1.4,1.5) {3};
    \node[tnodesm] (c) at (1.4,1.5) {10};
    \node[tnodesm] (d) at (-2.1,0.4) {1};
    \node[tnodesm] (e) at (-0.7,0.4) {6};
    \node[tnodesm] (f) at (1.4,0.4) {14};
    \node[tnodesm] (g) at (-1.1,-0.7) {4};
    \node[tnodesm] (h) at (-0.3,-0.7) {7};
    \node[tnodesm] (i) at (0.7,0.4) {13};
    \draw[tedge] (a) -- (b);  \draw[tedge] (a) -- (c);
    \draw[tedge] (b) -- (d);  \draw[tedge] (b) -- (e);
    \draw[tedge] (c) -- (f);  \draw[tedge] (e) -- (g);
    \draw[tedge] (e) -- (h);  \draw[tedge] (f) -- (i);
  \end{tikzpicture}
  \caption{A second tree for the traversal exercise.}
  \label{fig:check}
\end{figure}

\textbf{Solution.} Inorder visits left, node, right: 1, 3, 4, 6, 7, 8, 10,
13, 14. Preorder visits node, left, right: 8, 3, 1, 6, 4, 7, 10, 14, 13.
Postorder visits left, right, node: 1, 4, 7, 6, 3, 13, 14, 10, 8. Notice
that the inorder output is sorted, so this tree is also a BST; the check is
useful because an inorder that is not sorted proves the tree is not a BST.
\end{example}

\section{Expression Trees}

An arithmetic expression is a tree in disguise. Week~4 converted
\texttt{(a+b)*c} to postfix \texttt{a b + c *}. We can read that postfix left
to right with a stack of \emph{trees} and rebuild the expression:

\begin{itemize}
  \item An \textbf{operand} (a number or variable) becomes a one-node tree,
    which is pushed on the stack.
  \item An \textbf{operator} pops the \emph{right} tree, then the \emph{left}
    tree; it makes a new node holding the operator, with those two trees as
    children, and pushes the new tree.
  \item At the end, the stack holds exactly one tree: the expression.
\end{itemize}

The order of the two pops matters. For \texttt{a b +}, the first pop is
\texttt{b} (the right operand) and the second is \texttt{a} (the left
operand). Reversing them builds the tree for \texttt{b+a}, which happens to
have the same value but is not the same tree; for subtraction or division it
would be wrong.

\begin{example}
Build the expression tree for \texttt{a b + c *}, which is the postfix form
of \texttt{(a+b)*c}.

Read \texttt{a}: push leaf \texttt{a}. Read \texttt{b}: push leaf
\texttt{b}. Read \texttt{+}: pop \texttt{b} (right), pop \texttt{a} (left),
make a node \texttt{+} with those children, push it. Read \texttt{c}: push
leaf \texttt{c}. Read \texttt{*}: pop \texttt{c} (right), pop the
\texttt{+} tree (left), make the root \texttt{*} with those children. The
result is Figure~\ref{fig:expr}.
\end{example}

\begin{figure}[h]
  \centering
  \begin{tikzpicture}[every node/.style={font=\small}]
    \node[tnode] (star) at (0,2.4) {\texttt{*}};
    \node[tnode] (plus) at (-1.2,1.2) {\texttt{+}};
    \node[tnode] (c)    at (1.2,1.2)  {\texttt{c}};
    \node[tnode] (a)    at (-2.0,0.1) {\texttt{a}};
    \node[tnode] (b)    at (-0.4,0.1) {\texttt{b}};
    \draw[tedge] (star) -- (plus);
    \draw[tedge] (star) -- (c);
    \draw[tedge] (plus) -- (a);
    \draw[tedge] (plus) -- (b);
  \end{tikzpicture}
  \caption{The expression tree for \texttt{(a+b)*c}. Operators are internal
    nodes; operands are leaves. Cf.\ Karumanchi, Ch.~6, pp.224--368 (PDF)
    \cite{Karumanchi2017_data_structures_algorithms_made_easy}.}
  \label{fig:expr}
\end{figure}

To evaluate the tree, evaluate the left subtree, evaluate the right subtree,
then apply the operator --- exactly postorder. With $a=2$, $b=3$, $c=4$: leaf
\texttt{a} gives 2, leaf \texttt{b} gives 3, node \texttt{+} gives $2+3=5$,
leaf \texttt{c} gives 4, and the root \texttt{*} gives $5\times 4=20$. No
parentheses appear anywhere in the tree; the tree stores the grouping that
parentheses would otherwise have to spell out. That is why compilers build
trees.

Inorder on the same tree prints \texttt{a + b * c}, which is \emph{not} the
original expression: without parentheses, the usual precedence reads it as
\texttt{a + (b*c)}. To recover the meaning, put parentheses around every
operator's result as the recursion returns: first $(\texttt{a}+\texttt{b})$,
then $((\texttt{a}+\texttt{b})\times\texttt{c})$. The three walks match the
three notations from Week~4: inorder is infix (parentheses needed), preorder
is prefix, and postorder is postfix (and the evaluation order).

\section{Binary Search Trees}

\subsection{The ordering rule}

\begin{definitionbox}{binary search tree}
A \key{binary search tree} (BST) is a binary tree in which, for \emph{every}
node $x$, every key in the left subtree of $x$ is smaller than $x$'s key and
every key in the right subtree of $x$ is larger.
\end{definitionbox}

The rule must hold at every node, not only between a parent and its immediate
child. This is what makes search possible: at each node one comparison throws
away an entire subtree. Duplicates need a policy. The simplest is to reject a
duplicate; another is to keep a count in the node. Either is fine if it is
stated and applied consistently.

\begin{example}
Decide whether each tree in Figure~\ref{fig:isbst} is a BST.

\begin{figure}[h]
  \centering
  \begin{tikzpicture}[every node/.style={font=\scriptsize}]
    \node[tnodesm] (a) at (0,2.2) {5};
    \node[tnodesm] (b) at (-0.9,1.2) {3};
    \node[tnodesm] (c) at (0.9,1.2) {8};
    \node[tnodesm] (d) at (-1.4,0.2) {1};
    \node[tnodesm] (e) at (-0.4,0.2) {4};
    \node[tnodesm] (f) at (0.4,0.2) {7};
    \node[tnodesm] (g) at (1.4,0.2) {9};
    \draw[tedge] (a)--(b); \draw[tedge] (a)--(c);
    \draw[tedge] (b)--(d); \draw[tedge] (b)--(e);
    \draw[tedge] (c)--(f); \draw[tedge] (c)--(g);
    \node[text=positive, font=\scriptsize] at (0,-0.55) {\good{yes}};

    \node[tnodesm] (A) at (4.5,2.2) {5};
    \node[tnodesm] (B) at (3.6,1.2) {3};
    \node[tnodesm] (C) at (5.4,1.2) {8};
    \node[tnodesm] (D) at (3.1,0.2) {1};
    \node[tnodesmg] (E) at (4.1,0.2) {6};
    \node[tnodesm] (F) at (4.9,0.2) {7};
    \node[tnodesm] (G) at (5.9,0.2) {9};
    \draw[tedge] (A)--(B); \draw[tedge] (A)--(C);
    \draw[tedge] (B)--(D); \draw[tedge] (B)--(E);
    \draw[tedge] (C)--(F); \draw[tedge] (C)--(G);
    \node[text=negative, font=\scriptsize] at (4.5,-0.55)
      {\bad{no: 6 $>$ 5 but lies left of 5}};
  \end{tikzpicture}
  \caption{Testing the BST rule at every node, not only at parent--child
    pairs.}
  \label{fig:isbst}
\end{figure}

\textbf{Solution.} The left tree is a BST. The right tree is not: 6 is a
child of 3 and $6>3$ looks locally fine, but 6 lies in the left subtree of
the root 5, and $6>5$. The violation is against the root, not against the
parent. Checking only parent--child pairs is not enough.
\end{example}

\subsection{Search}

Search starts at the root and compares. If the target equals the node's key,
stop. If it is smaller, go left; if it is larger, go right. If the walk
reaches \texttt{NULL}, the key is absent.

\begin{figure}[h]
  \centering
  \begin{tikzpicture}[every node/.style={font=\small}]
    \node[tnode] (n50) at (0,2.8)   {50};
    \node[tnode] (n30) at (-1.7,1.7) {30};
    \node[tnode] (n70) at (1.7,1.7)  {70};
    \node[tnode] (n20) at (-2.5,0.6) {20};
    \node[tnodeg] (n40) at (-0.9,0.6) {40};
    \node[tnode] (n60) at (0.9,0.6)  {60};
    \node[tnode] (n80) at (2.5,0.6)  {80};
    \draw[tedge] (n50) -- (n30);
    \draw[tedge] (n50) -- (n70);
    \draw[tedge] (n30) -- (n20);
    \draw[tedge] (n30) -- (n40);
    \draw[tedge] (n70) -- (n60);
    \draw[tedge] (n70) -- (n80);
    \draw[-{Latex[length=2mm]}, primary-gold, thick]
      (n50) to[bend left=12] (n30);
    \draw[-{Latex[length=2mm]}, primary-gold, thick]
      (n30) to[bend left=12] (n40);
  \end{tikzpicture}
  \caption{Searching for 40: $40<50$ go left, $40>30$ go right, found.}
  \label{fig:search}
\end{figure}

\begin{example}
Trace a search for 40 in Figure~\ref{fig:search} and count its comparisons.

Compare 40 with 50: smaller, go left. Compare with 30: larger, go right.
Compare with 40: equal, found. Three comparisons, not seven. The cost is the
number of nodes on the path, which is at most height $+1$.
\end{example}

In C the loop is iterative, so it uses $O(1)$ extra space:

\begin{verbatim}
struct node *bst_search(struct node *t, int key) {
    while (t != NULL && t->data != key) {
        if (key < t->data) t = t->left;
        else               t = t->right;
    }
    return t;          /* NULL means "not present" */
}
\end{verbatim}

The worst-case time is $O(h)$. A balanced tree has $h=O(\log n)$; a chain has
$h=n-1$. This is the same loop as binary search on a sorted array, with
\texttt{left}/\texttt{right} pointers instead of index arithmetic.

\subsection{Insertion}

Insertion searches for the key. If the key is already present, reject it (or
count it). If the search falls off the tree at \texttt{NULL}, the new node
goes there. A new key always becomes a leaf.

\begin{example}
Insert $50,30,70,20,40,60,80$ into an empty BST.

50 becomes the root. 30 is less than 50, so it becomes 50's left child. 70 is
greater than 50, so it becomes 50's right child. 20 is less than 50 and less
than 30, so it becomes 30's left child. 40 is less than 50 but greater than
30, so it becomes 30's right child. 60 is greater than 50 but less than 70,
so it becomes 70's left child. 80 is greater than 50 and greater than 70, so
it becomes 70's right child. The result is Figure~\ref{fig:anatomy}.
\end{example}

Insertion never rearranges the existing nodes. That keeps it simple, but it
is also the weakness that Section~11.6 fixes: the shape is whatever the
arrival order produces.

\subsection{Minimum and maximum}

The minimum is found by starting at the root and taking \texttt{left} until
there is no left child. In Figure~\ref{fig:anatomy}: $50\to30\to20$, so the
minimum is 20. The maximum is the mirror: $50\to70\to80$, so the maximum is
80. Both are $O(h)$, one root-to-leaf walk. They are the first ordered
operations a hash table cannot do at all.

\subsection{Deletion}

Deletion has three cases.

\textbf{Case 1 --- leaf.} Remove the node and set the parent's pointer to
\texttt{NULL}. Deleting 20 from Figure~\ref{fig:anatomy} makes 30's left
pointer \texttt{NULL}.

\textbf{Case 2 --- one child.} Splice the node out: the parent points
straight to the node's only child. If 30 has only the right child 40,
deleting 30 makes 40 the left child of 50. Both cases preserve the BST rule,
because the subtree that replaces the deleted node was already on the correct
side.

\textbf{Case 3 --- two children.} The node cannot simply be removed, because
two subtrees need a parent. Find the \key{successor}: the smallest key in the
right subtree, obtained by going right once and then left as far as possible.
Copy the successor's key into the node, then delete the successor from the
right subtree. The successor has no left child, so its deletion is Case 1 or
Case 2.

\begin{example}
Delete 50 from Figure~\ref{fig:anatomy}.

50 has two children. Its right subtree is rooted at 70. Go right to 70, then
left to 60: the successor is 60. Copy 60 into the root, then delete the old
60 node. The old 60 was a leaf, so Case 1 removes it. The result is
Figure~\ref{fig:delete}: 60 is the new root, with 30 on the left (and 20, 40
below it) and 70 on the right (with 80 below it).
\end{example}

\begin{figure}[h]
  \centering
  \begin{tikzpicture}[every node/.style={font=\scriptsize}]
    \foreach \x/\y/\v in {1.2/2.2/50, 0.45/1.3/30, 1.95/1.3/70,
                          0.05/0.4/20, 0.85/0.4/40, 1.55/0.4/60, 2.35/0.4/80}
      \node[tnodesm] at (\x,\y) {\v};
    \draw[tedge] (1.2,2.2)--(0.45,1.3); \draw[tedge] (1.2,2.2)--(1.95,1.3);
    \draw[tedge] (0.45,1.3)--(0.05,0.4); \draw[tedge] (0.45,1.3)--(0.85,0.4);
    \draw[tedge] (1.95,1.3)--(1.55,0.4); \draw[tedge] (1.95,1.3)--(2.35,0.4);
    \node[tnodesmg] at (1.55,0.4) {60};
    \node[text=neutral] at (1.2,-0.35) {before: delete 50};

    \foreach \x/\y/\v in {6.0/2.2/60, 5.25/1.3/30, 6.75/1.3/70,
                          4.85/0.4/20, 5.65/0.4/40, 7.15/0.4/80}
      \node[tnodesm] at (\x,\y) {\v};
    \draw[tedge] (6.0,2.2)--(5.25,1.3); \draw[tedge] (6.0,2.2)--(6.75,1.3);
    \draw[tedge] (5.25,1.3)--(4.85,0.4); \draw[tedge] (5.25,1.3)--(5.65,0.4);
    \draw[tedge] (6.75,1.3)--(7.15,0.4);
    \node[tnodesmg] at (6.0,2.2) {60};
    \node[text=neutral] at (6.0,-0.35) {after: 60 replaces 50};
  \end{tikzpicture}
  \caption{Deleting a node with two children: the successor 60 replaces 50,
    and the old 60 node is removed from the right subtree. Cf.\ Wirth,
    Sec.~4.4.4, p.139 \cite{Wirth2004_algorithms_data_structures}.}
  \label{fig:delete}
\end{figure}

The recursive C version follows the same three cases:

\begin{verbatim}
struct node *bst_delete(struct node *t, int key) {
    if (t == NULL) return NULL;
    if (key < t->data)      t->left  = bst_delete(t->left,  key);
    else if (key > t->data) t->right = bst_delete(t->right, key);
    else {                                  /* found it */
        if (t->left == NULL)  { struct node *r = t->right; free(t); return r; }
        if (t->right == NULL) { struct node *l = t->left;  free(t); return l; }
        struct node *s = t->right;           /* successor */
        while (s->left != NULL) s = s->left;
        t->data  = s->data;
        t->right = bst_delete(t->right, s->data);
    }
    return t;
}
\end{verbatim}

The recursive call in the two-child case removes the successor, which has no
left child. Every node is visited once, so the cost is $O(h)$.

\subsection{The cost is the height}

Every BST operation is one root-to-leaf walk, so search, insert and delete
all cost $O(h)$. Table~\ref{tab:bstcost} gives the three shapes that matter.

\begin{table}[h]
  \centering
  {\footnotesize
  \begin{tabular}{l c c}
    \toprule
    \textbf{shape} & \textbf{height $h$} & \textbf{search / insert / delete} \\
    \midrule
    balanced (best) & $\lfloor \log_2 n\rfloor$ & $O(\log n)$ \\
    random insertions (average) & $\approx 1.39\log_2 n$ & $O(\log n)$ expected \\
    chain (worst) & $n-1$ & $O(n)$ \\
    \bottomrule
  \end{tabular}}
  \caption{The cost of a BST is its height. The average row assumes random
    arrival order.}
  \label{tab:bstcost}
\end{table}

The average row assumes keys arrive in random order. Real data is often
sorted --- names, timestamps, roll numbers --- and that is the dangerous case.

\begin{example}
Insert $10,20,30,40,50$ into an empty BST, in that order.

Every new key is larger than all existing keys, so every insert goes right.
The result is a right chain of height 4 (Figure~\ref{fig:chain}, left). A
search for 50 visits $10,20,30,40,50$: five comparisons, exactly what a
linked list would cost. The same five keys can also be arranged as a
balanced tree of height 2 (Figure~\ref{fig:chain}, right), where the same
search takes three comparisons. Nothing in the BST insert rule prevents the
chain; we need a rule that repairs the shape.
\end{example}

\begin{figure}[h]
  \centering
  \begin{tikzpicture}[every node/.style={font=\scriptsize}]
    \foreach \y/\v in {3.2/10, 2.4/20, 1.6/30, 0.8/40, 0/50}
      \node[tnodesm] at (0,\y) {\v};
    \foreach \y in {2.8,2.0,1.2,0.4}
      \draw[tedge] (0,\y+0.4) -- (0,\y-0.4);
    \node[text=negative, align=center] at (0,-0.65)
      {\bad{height 4}\\ search(50) $=5$ comparisons};

    \foreach \x/\y/\v in {4.5/2.4/30, 3.8/1.4/20, 5.2/1.4/40,
                          3.45/0.4/10, 4.15/0.4/25, 4.85/0.4/35, 5.55/0.4/50}
      \node[tnodesm] at (\x,\y) {\v};
    \draw[tedge] (4.5,2.4)--(3.8,1.4); \draw[tedge] (4.5,2.4)--(5.2,1.4);
    \draw[tedge] (3.8,1.4)--(3.45,0.4); \draw[tedge] (3.8,1.4)--(4.15,0.4);
    \draw[tedge] (5.2,1.4)--(4.85,0.4); \draw[tedge] (5.2,1.4)--(5.55,0.4);
    \node[text=positive, align=center] at (4.5,-0.65)
      {\good{height 2}\\ search(50) $=3$ comparisons};
  \end{tikzpicture}
  \caption{The same five keys in two shapes. The chain on the left costs as
    much as a linked list; the balanced tree on the right does not.}
  \label{fig:chain}
\end{figure}

\section{AVL Trees}

\subsection{The balance rule and the balance factor}

\begin{definitionbox}{AVL tree}
An \key{AVL tree} is a BST in which, for every node, the heights of the left
and right subtrees differ by at most $1$.
\end{definitionbox}

The \key{balance factor} of a node $x$ is
\[
  \mathrm{bf}(x) = \mathrm{height}(\text{left child of }x)
                 - \mathrm{height}(\text{right child of }x).
\]
The AVL rule is $\mathrm{bf}(x)\in\{-1,0,+1\}$ for every node. If an insert
or delete makes some $|\mathrm{bf}|=2$, the tree is repaired with a
\key{rotation}: a local pointer change that lowers the tall side. The tree is
named after Adelson-Velsky and Landis (1962). We keep the height convention
from Section~11.2: a leaf has height $0$, the empty tree has height $-1$.

\begin{example}
Compute the balance factor of every node in Figure~\ref{fig:bf}.

\begin{figure}[h]
  \centering
  \begin{tikzpicture}[every node/.style={font=\small}]
    \node[tnode] (n50) at (0,2.5) {50};
    \node[tnode] (n30) at (-1.3,1.4) {30};
    \node[tnode] (n70) at (1.3,1.4) {70};
    \node[tnode] (n20) at (-1.9,0.3) {20};
    \draw[tedge] (n50)--(n30); \draw[tedge] (n50)--(n70);
    \draw[tedge] (n30)--(n20);
    \node[text=primary-gold, font=\scriptsize] at (0.0,3.38) {bf $=+1$};
    \node[text=primary-gold, font=\scriptsize] at (-1.3,2.30) {bf $=+1$};
    \node[text=primary-gold, font=\scriptsize] at (1.3,2.30) {bf $=0$};
    \node[text=primary-gold, font=\scriptsize] at (-1.9,-0.45) {bf $=0$};
  \end{tikzpicture}
  \caption{Balance factors on a valid AVL tree.}
  \label{fig:bf}
\end{figure}

\textbf{Solution.} Node 20 is a leaf, so its height is 0 and
$\mathrm{bf}(20)=0$. Node 70 is a leaf, so $\mathrm{bf}(70)=0$. Node 30 has
left child 20 (height 0) and no right child (height $-1$), so
$\mathrm{bf}(30)=0-(-1)=+1$. Node 50 has left subtree 30 (height 1) and
right subtree 70 (height 0), so $\mathrm{bf}(50)=1-0=+1$. Every factor lies in
$\{-1,0,+1\}$, so this is an AVL tree.
\end{example}

\subsection{A rotation is three pointer moves}

A rotation lifts the middle key and pushes the old root down. Figure~\ref{fig:rot}
shows a right rotation at 30. The node 20 rises; 30 becomes its right child;
20's old right subtree \texttt{T2} becomes 30's left subtree. The inorder
sequence is $10,20,\texttt{T2},30$ both before and after, so the BST rule is
preserved. A rotation rearranges the shape, never the order.

\begin{figure}[h]
  \centering
  \begin{tikzpicture}[every node/.style={font=\scriptsize}]
    \node[tnodesm] (y) at (1,2.2) {30};
    \node[tnodesm] (x) at (0,1.2) {20};
    \node[tnodesm] (t1) at (-0.5,0.2) {10};
    \node[draw=neutral, fill=white, minimum width=0.55cm,
      minimum height=0.45cm] (t2) at (0.5,0.2) {T2};
    \draw[tedge] (y)--(x); \draw[tedge] (x)--(t1); \draw[tedge] (x)--(t2);
    \node[text=neutral] at (0.5,-0.55) {before};

    \draw[-{Latex[length=2.5mm]}, primary-gold, thick]
      (2.0,1.2) -- (2.9,1.2);

    \node[tnodesm] (X) at (5.0,2.2) {20};
    \node[tnodesm] (T1) at (4.35,1.2) {10};
    \node[tnodesm] (Y) at (5.65,1.2) {30};
    \node[draw=neutral, fill=white, minimum width=0.55cm,
      minimum height=0.45cm] (T2b) at (5.0,0.2) {T2};
    \draw[tedge] (X)--(T1); \draw[tedge] (X)--(Y); \draw[tedge] (Y)--(T2b);
    \node[text=neutral] at (5.0,-0.55) {after right rotation at 30};
  \end{tikzpicture}
  \caption{A right rotation at 30. The middle key 20 rises; the old root 30
    becomes its right child; the middle key's old right subtree moves to the
    old root's left. Cf.\ Wirth, Sec.~4.5, pp.142--149
    \cite{Wirth2004_algorithms_data_structures}.}
  \label{fig:rot}
\end{figure}

\subsection{The four cases}

There are four ways a single insertion can unbalance the lowest bad node,
named by the path from that node to the inserted key.

\begin{center}
{\footnotesize
\begin{tabular}{l l l}
  \toprule
  \textbf{case} & \textbf{path from the unbalanced node} & \textbf{repair} \\
  \midrule
  LL & left child, then left child  & right rotation once \\
  RR & right child, then right child & left rotation once \\
  LR & left child, then right child & left rotate child, then right rotate node \\
  RL & right child, then left child & right rotate child, then left rotate node \\
  \bottomrule
\end{tabular}}
\end{center}

Read the name as the \emph{path}, not as the first rotation. LR means ``went
left, then right''. The first rotation turns LR into LL; the second rotation
finishes the job. Only the nodes on the path from the inserted leaf back to
the root can change balance, and the lowest violation is repaired first.

\begin{example}
Trace the LL and RR repairs.

\textbf{LL.} Insert $30,20,10$ in that order. The path from 30 is left
(20), then left (10), so this is LL. Node 30 has
$\mathrm{bf}=+2$: left subtree height 1, right subtree height $-1$. A right
rotation at 30 makes 20 the root, with 10 on the left and 30 on the right.
The three-node tree now has height 1 and every factor is 0.

\textbf{RR.} Insert $10,20,30$. The path is right, then right, so this is
RR. A left rotation at 10 makes 20 the root, with 10 left and 30 right. LL
and RR are mirror images; one rotation each.
\end{example}

\begin{figure}[h]
  \centering
  \begin{tikzpicture}[every node/.style={font=\scriptsize}]
    % LL: 30-20-10 -> right rotation at 30 -> 20 root.
    \node[tnodesm] (l1) at (0,2.2) {30};
    \node[tnodesm] (l2) at (0,1.2) {20};
    \node[tnodesm] (l3) at (0,0.2) {10};
    \draw[tedge] (l1)--(l2); \draw[tedge] (l2)--(l3);
    \node[text=neutral, font=\tiny] at (-0.15,-0.45) {LL: 30--20--10};
    \draw[-{Latex[length=2mm]}, primary-gold, thick] (0.55,1.2)--(1.15,1.2);
    \node[tnodesm] (L1) at (1.8,2.2) {20};
    \node[tnodesm] (L2) at (1.3,1.2) {10};
    \node[tnodesm] (L3) at (2.3,1.2) {30};
    \draw[tedge] (L1)--(L2); \draw[tedge] (L1)--(L3);
    \node[text=positive, font=\tiny] at (1.95,-0.45) {\good{right rotate}};

    % RR: 10-20-30 -> left rotation at 10 -> 20 root.
    \node[tnodesm] (r1) at (4.6,2.2) {10};
    \node[tnodesm] (r2) at (4.6,1.2) {20};
    \node[tnodesm] (r3) at (4.6,0.2) {30};
    \draw[tedge] (r1)--(r2); \draw[tedge] (r2)--(r3);
    \node[text=neutral, font=\tiny] at (4.45,-0.45) {RR: 10--20--30};
    \draw[-{Latex[length=2mm]}, primary-gold, thick] (5.15,1.2)--(5.75,1.2);
    \node[tnodesm] (R1) at (6.4,2.2) {20};
    \node[tnodesm] (R2) at (5.9,1.2) {10};
    \node[tnodesm] (R3) at (6.9,1.2) {30};
    \draw[tedge] (R1)--(R2); \draw[tedge] (R1)--(R3);
    \node[text=positive, font=\tiny] at (6.55,-0.45) {\good{left rotate}};
  \end{tikzpicture}
  \caption{LL and RR each need one rotation. The middle key rises; the old
    root becomes a child.}
  \label{fig:llrr}
\end{figure}

\begin{example}
Trace the LR and RL double rotations.

\textbf{LR.} Insert $30,10,20$. The path from 30 is left (10), then right
(20). Left-rotate 10 around 20 first: 30's left child becomes 20, and 20's
left child becomes 10. The path is now LL. Right-rotate 30 around 20: 20
becomes the root, with 10 left and 30 right.

\textbf{RL.} Insert $10,30,20$. The path from 10 is right (30), then left
(20). Right-rotate 30 around 20: 10's right child becomes 20, and 20's right
child becomes 30. The path is now RR. Left-rotate 10 around 20: 20 becomes
the root, with 10 left and 30 right. In both cases the inorder sequence stays
$10,20,30$.
\end{example}

\begin{figure}[h]
  \centering
  \begin{tikzpicture}[every node/.style={font=\scriptsize}]
    \node[tnodesm] (a1) at (0,2.2) {30};
    \node[tnodesm] (a2) at (-0.7,1.2) {10};
    \node[tnodesm] (a3) at (0.0,0.2) {20};
    \draw[tedge] (a1)--(a2); \draw[tedge] (a2)--(a3);
    \node[text=neutral, font=\tiny] at (0,-0.45) {start};

    \draw[-{Latex[length=2mm]}, primary-gold, thick] (0.6,1.2)--(1.2,1.2);
    \node[tnodesm] (b1) at (2.1,2.2) {30};
    \node[tnodesm] (b2) at (1.5,1.2) {20};
    \node[tnodesm] (b3) at (1.0,0.2) {10};
    \draw[tedge] (b1)--(b2); \draw[tedge] (b2)--(b3);
    \node[text=neutral, font=\tiny] at (2.1,-0.45) {left rotate 10};

    \draw[-{Latex[length=2mm]}, primary-gold, thick] (2.7,1.2)--(3.3,1.2);
    \node[tnodesm] (c1) at (4.2,2.2) {20};
    \node[tnodesm] (c2) at (3.7,1.2) {10};
    \node[tnodesm] (c3) at (4.7,1.2) {30};
    \draw[tedge] (c1)--(c2); \draw[tedge] (c1)--(c3);
    \node[text=positive, font=\tiny] at (4.2,-0.45) {\good{right rotate 30}};
  \end{tikzpicture}
  \caption{The LR case: left-rotate the left child, then right-rotate the
    node. The intermediate tree is an LL case.}
  \label{fig:lr}
\end{figure}

\begin{figure}[h]
  \centering
  \begin{tikzpicture}[every node/.style={font=\scriptsize}]
    \node[tnodesm] (a1) at (0,2.2) {10};
    \node[tnodesm] (a2) at (0.7,1.2) {30};
    \node[tnodesm] (a3) at (0.0,0.2) {20};
    \draw[tedge] (a1)--(a2); \draw[tedge] (a2)--(a3);
    \node[text=neutral, font=\tiny] at (0,-0.45) {start};

    \draw[-{Latex[length=2mm]}, primary-gold, thick] (0.6,1.2)--(1.2,1.2);
    \node[tnodesm] (b1) at (2.1,2.2) {10};
    \node[tnodesm] (b2) at (2.7,1.2) {20};
    \node[tnodesm] (b3) at (3.2,0.2) {30};
    \draw[tedge] (b1)--(b2); \draw[tedge] (b2)--(b3);
    \node[text=neutral, font=\tiny] at (2.1,-0.45) {right rotate 30};

    \draw[-{Latex[length=2mm]}, primary-gold, thick] (2.7,1.2)--(3.3,1.2);
    \node[tnodesm] (c1) at (4.2,2.2) {20};
    \node[tnodesm] (c2) at (3.7,1.2) {10};
    \node[tnodesm] (c3) at (4.7,1.2) {30};
    \draw[tedge] (c1)--(c2); \draw[tedge] (c1)--(c3);
    \node[text=positive, font=\tiny] at (4.2,-0.45) {\good{left rotate 10}};
  \end{tikzpicture}
  \caption{The RL case: right-rotate the right child, then left-rotate the
    node. The intermediate tree is an RR case.}
  \label{fig:rl}
\end{figure}

\begin{example}
Insert $30,20,10,25,28$ into an empty AVL tree, repairing after each
insertion.

\textbf{30, 20, 10.} The path from 30 is LL. Right-rotate 30: 20 becomes the
root, with 10 left and 30 right.

\textbf{25.} It goes left of 30. The factors are
$\mathrm{bf}(30)=+1$ and $\mathrm{bf}(20)=-1$, both legal; no rotation.

\textbf{28.} It goes right of 25. Now 30's left subtree has height 1 (25
with 28) and its right subtree is empty (height $-1$), so
$\mathrm{bf}(30)=+2$. The path from 30 is left (25), then right (28): an LR
case. Left-rotate 25 around 28, then right-rotate 30 around 28. The subtree
root becomes 28, with 25 left and 30 right. The whole tree is 20 at the root,
10 on the left, and 28 on the right (Figure~\ref{fig:avlfinal}); its height
is 2 and every factor is legal.
\end{example}

\begin{figure}[h]
  \centering
  \begin{tikzpicture}[every node/.style={font=\scriptsize}]
    \node[tnodesm] (r) at (0,2.4) {20};
    \node[tnodesm] (l) at (-1.0,1.3) {10};
    \node[tnodesm] (m) at (1.0,1.3) {28};
    \node[tnodesm] (ml) at (0.5,0.2) {25};
    \node[tnodesm] (mr) at (1.5,0.2) {30};
    \draw[tedge] (r)--(l); \draw[tedge] (r)--(m);
    \draw[tedge] (m)--(ml); \draw[tedge] (m)--(mr);
  \end{tikzpicture}
  \caption{The AVL tree after inserting $30,20,10,25,28$. Height 2.}
  \label{fig:avlfinal}
\end{figure}

\subsection{Insertion in C}

The node stores its height so the balance check is $O(1)$ per node:

\begin{verbatim}
int height(struct node *t) { return t ? t->height : -1; }
int bf(struct node *t)     { return t ? height(t->left) - height(t->right) : 0; }

struct node *rotate_right(struct node *y) {
    struct node *x  = y->left;          /* x rises */
    struct node *t2 = x->right;         /* T2 changes parent */
    x->right = y;  y->left = t2;
    y->height = 1 + max(height(y->left),  height(y->right));
    x->height = 1 + max(height(x->left),  height(x->right));
    return x;                            /* new subtree root */
}
/* rotate_left is the mirror image. */
\end{verbatim}

After a normal BST insert, walk back up. At each node update \texttt{height},
compute \texttt{bf}, and apply the LL/RR/LR/RL case whose factor is $\pm2$.

\subsection{Why the height stays logarithmic}

Let $N(h)$ be the \emph{fewest} nodes an AVL tree of height $h$ can have. The
worst shape has one subtree of height $h-1$ and the other of height $h-2$:
\[
  N(h) = N(h-1) + N(h-2) + 1,\qquad N(0)=1,\ N(1)=2.
\]
This is the Fibonacci recurrence (shifted), so $N(h)$ grows exponentially:
$1,2,4,7,12,20,33,\dots$ Therefore $h=O(\log n)$, with the precise bound
\[
  h \le 1.44\log_2(n+2) - 0.33.
\]
For $n=1{,}000{,}000$ this gives $h\le 28$, against $999{,}999$ for the chain.
Search, insert and delete are all $O(\log n)$ worst case. You do not need the
constant $1.44$ in an exam; you do need the recurrence and the conclusion.

\subsection{Deletion}

Delete exactly as in a BST (the three cases), then walk from the parent of the
removed node back up to the root. At each node recompute the height and
balance factor; if $|\mathrm{bf}|=2$, apply the same four cases. Deletion can
need \emph{several} rotations, not just one: removing a node can unbalance
several ancestors. Insertion needs at most one repair on the path. Deletion is
still $O(\log n)$, because the walk has $\log n$ levels and each rotation is
$O(1)$.

The invariant to remember is that a rotation preserves the BST rule and the
inorder sequence. That is the only reason we are allowed to repair the shape
without touching the data.

\section{Choosing Between a Tree, a Hash Table, and a Sorted Array}

\begin{center}
{\footnotesize
\begin{tabular}{l c c c c}
  \toprule
  \textbf{structure} & \textbf{search} & \textbf{insert} & \textbf{ordered ops} & \textbf{space} \\
  \midrule
  sorted array & $O(\log n)$ & $O(n)$ & yes & $O(n)$ \\
  hash table   & $O(1)$ expected & $O(1)$ expected & no & $O(n)$ \\
  plain BST    & $O(h)$ & $O(h)$ & yes & $O(n)$ \\
  AVL tree     & $O(\log n)$ & $O(\log n)$ & yes & $O(n)$ \\
  \bottomrule
\end{tabular}}
\end{center}

The choice follows the workload, not the fashion. If the only question is
``is this key present?'' and order is irrelevant, a hash table is the right
tool. If the questions include sorted output, next/previous, range queries,
minimum or maximum, use an AVL tree. If the data is static, read many times
and never updated, a sorted array is simplest. ``Trees are better than
hashing'' is false; ``trees are better when order is part of the question'' is
true.

\section{Unit III in One Page}

\begin{center}
{\footnotesize
\begin{tabular}{c l l l}
  \toprule
  \textbf{week} & \textbf{idea} & \textbf{key cost} & \textbf{key trap} \\
  \midrule
  8 & search; elementary sorts &
    $\log n$ search; $n^2$ sorts & binary search needs sorted data \\
  9 & efficient sorts &
    $n\log n$; radix $d(n+b)$ & quicksort is $n^2$ on sorted input \\
  10 & hashing &
    $O(1)$ expected lookup & order is lost; tombstone on delete \\
  11 & trees, BST, AVL &
    $O(\log n)$ if balanced & a plain BST can become a chain \\
  \bottomrule
\end{tabular}}
\end{center}

The unit's single theme is that the cost of an operation is decided by the
shape of the data, not by the name of the structure. Week~12 removes the last
restriction: trees have one root and no cycles; graphs may have many paths,
cycles, and no root at all.

\section{Three Classic Mistakes and a Summary}

\textbf{Mistake 1: ``It is a binary tree, so search is $O(\log n)$.''} Only a
\emph{short} BST is logarithmic. Inserting sorted keys builds a chain, and
search becomes $O(n)$.

\textbf{Mistake 2: ``Inorder is always sorted.''} Only for a BST. On an
arbitrary binary tree, inorder is just left--node--right.

\textbf{Mistake 3: ``A rotation changes the keys.''} A rotation changes
pointers and heights, never the inorder sequence. If the inorder changes, the
rotation is wrong.

The week in five lines:

\begin{itemize}
  \item \key{Tree vocabulary} --- root, leaf, parent, child, sibling, subtree;
    depth counts down from the root, height counts up from a leaf, and height
    controls the cost.
  \item \key{Traversals} --- inorder (left, node, right), preorder (node,
    left, right), postorder (left, right, node), and level order (queue).
    Each is $O(n)$; each has a job.
  \item \key{Expression trees} --- operands are leaves, operators are
    internal nodes; postorder evaluates, preorder gives prefix, inorder gives
    infix (with parentheses).
  \item \key{BST} --- left subtree $<$ node $<$ right subtree at every node.
    Search, insert, min, max and delete are $O(h)$; delete has three cases,
    and the two-child case uses the successor.
  \item \key{AVL} --- keep every balance factor in $\{-1,0,+1\}$; repair with
    LL/RR (one rotation) or LR/RL (two). Height, and therefore every
    operation, is $O(\log n)$.
\end{itemize}

Next week: graphs --- many connections, cycles, and no single root. BFS and
DFS, Prim and Kruskal, and Dijkstra's shortest paths.

\subsection*{Three Quick Checks}

\begin{enumerate}
  \item Insert $5,3,8,1,4,7,9$ into an empty BST. What is the inorder output,
    and what is the height? \emph{Answer:} $1,3,4,5,7,8,9$; height 2.
  \item Delete 5 from that tree. Which key replaces it, and why?
    \emph{Answer:} 7, the smallest key in the right subtree (the successor).
  \item Insert $10,20,30$ into an empty AVL tree. Draw the result and name
    the rotation. \emph{Answer:} the RR case; one left rotation at 10 makes
    20 the root, with 10 left and 30 right.
\end{enumerate}

\subsection*{Exercises}

\begin{enumerate}
  \item Insert $7,4,9,2,6,8,10,3,5$ into an empty BST in that order, then
    give its inorder, preorder, postorder and level-order traversals. What is
    its height? Is it an AVL tree?
  \item Insert $50,25,75,10,35,60,90,30,40$ into an empty BST. Give the
    inorder output and the height. Then delete 25 and, separately, delete 50
    from the original tree, showing the successor used in each case.
  \item Insert $40,20,60,10,30,25$ into an empty AVL tree. Name the rotation
    (or rotations) needed and draw the final tree.
  \item Compute $N(0)$ through $N(6)$ for the minimum-node AVL recurrence,
    and draw a 7-node AVL tree of height 3.
  \item Build the expression tree for the postfix string
    \texttt{5 3 + 8 2 - *}. Evaluate it, and give its prefix and fully
    parenthesised infix forms.
\end{enumerate}

\subsection*{Solutions}

\begingroup\small
\begin{enumerate}
  \item The BST is 7 at the root; 4 left (2 with right 3, 6 with left 5);
    9 right (8, 10). Inorder: 2, 3, 4, 5, 6, 7, 8, 9, 10. Preorder: 7, 4, 2,
    3, 6, 5, 9, 8, 10. Postorder: 3, 2, 5, 6, 4, 8, 10, 9, 7. Level order:
    7, 4, 9, 2, 6, 8, 10, 3, 5. Height 3 (the longest path is
    $7\to4\to2\to3$). It \emph{is} AVL: the left subtree of 7 has height 2
    and the right has height 1, so $\mathrm{bf}(7)=+1$; node 4 has equal
    subtree heights and factor 0, and nodes 2 and 6 have factors $-1$ and
    $+1$. The extra level on the left is allowed because the right side is
    only one level shorter.
  \item Inorder: 10, 25, 30, 35, 40, 50, 60, 75, 90. Height 3 (the path
    $50\to25\to35\to30$). Deleting 25: its right subtree is rooted at 35, and
    the leftmost node there is 30, so 30 replaces 25 and the old 30 leaf is
    removed. Deleting 50: its right subtree is rooted at 75, and the
    leftmost node there is 60, so 60 replaces 50 and the old 60 leaf is
    removed; the right side becomes 75 with right child 90.
  \item Inserting $40,20,60,10,30,25$: after the first five the tree is 40 at
    the root, 20 left (10, 30) and 60 right. Inserting 25 makes it the left
    child of 30. Now $\mathrm{bf}(40)=+2$ (left subtree height 2, right
    subtree height 0) and the path from 40 is left (20) then right (30): an
    LR case. Left-rotate 20 around 30, then right-rotate 40 around 30. Final
    tree: 30 at the root, 20 left (with 10 and 25), 40 right (with 60 as its
    right child). Every factor is in $\{-1,0,+1\}$ and the height is 2.
  \item $N(0)=1$, $N(1)=2$, $N(2)=4$, $N(3)=7$, $N(4)=12$, $N(5)=20$,
    $N(6)=33$. A 7-node AVL tree of height 3: root 10; left child 5 with
    left 2 (left child 1) and right 7; right child 15 with left 12. The root
    has subtree heights 2 and 1, so $\mathrm{bf}(10)=+1$; node 5 has heights
    1 and 0, so $\mathrm{bf}(5)=+1$; node 2 has heights 1 and $-1$, so
    $\mathrm{bf}(2)=+1$. All factors are legal and the height is 3.
  \item Read \texttt{5} (push), \texttt{3} (push), \texttt{+} (pop 3 right,
    5 left, push \texttt{+}); read \texttt{8} (push), \texttt{2} (push),
    \texttt{-} (pop 2 right, 8 left, push \texttt{-}); read \texttt{*}
    (pop the \texttt{-} tree right, the \texttt{+} tree left, push
    \texttt{*}). The root is \texttt{*}, with left \texttt{+}(5,3) and right
    \texttt{-}(8,2). Evaluate: $5+3=8$, $8-2=6$, $8\times6=48$. Prefix:
    \texttt{* + 5 3 - 8 2}. Fully parenthesised infix:
    $((5+3)\times(8-2))$.
\end{enumerate}
\endgroup

\bibliographystyle{plain}
\bibliography{../../Bibliography_base}

\end{document}
